\section{Saturate eight signed lanes}\label{pcm-vector}
Load the two groups as byte elements into vector registers. Their paired low/high
bytes form eight halfword lanes in the pinned little-endian environment. The same
bits can then be interpreted by halfword arithmetic without moving them to new
storage. \texttt{adds\_s.h} adds signed halfword values and clamps each lane to its
signed range. The unsigned saturation instruction implements a different operation.

Also compute \texttt{addv.h}, which retains each sum's low halfword bits. Comparing
the saturated and wrapping vectors yields equal lanes exactly for the samples that
did not clip. An out-of-range positive sum wraps to a negative value while saturating
to 32767. An out-of-range negative sum wraps to a nonnegative value while saturating
to minus 32768. Thus no clipped lane can accidentally compare equal in this two-input
domain.

\texttt{ceq.h} represents true with all sixteen bits set and false with all bits
clear. Invert that equality result with \texttt{nor.v} to obtain all-one lanes for
clipped samples and all-zero lanes for in-range samples. This representation is a
bit mask, not the integer number one per clipped sample.

A group store can use the saturated vector after both inputs have been loaded and
the counter evidence has been retained. The guided sequence completes the counter
first and then stores. Equivalent implementations may order these internal operations
differently while preserving inputs needed for the allowed aliases and the exact
public result.
